% !TEX root = ../main.tex

\section{Original Contribution}
\label{sec:simulations}


%TODO: What is original Contribution?
% Your contribution to this project? What have you done from everything you have learned.  

\subsection*{Stability Analysis}

% Step 1: Tell the reader what I'm going to do in this section
The Gray-Scott system is a nonlinear dynamical system whose equations cannot be solved analytically. Therefore, we consider a method known as stability analysis which will predict the behaviour of the system and precisely identify the set of values for the parameters $F$ and $k$ that lead to the turing patterns we hinted in the previous section. The paramteter values are found using bifurcation analysis

Steps I'll be doing for stability analysis:
\begin{itemize}
	\item  Firstly, I will drop the diffusion terms to show that the system is uniformly stable.
	\item But before we do that, I solve the system for it's equillibrium points.
	\item Then classify the system by substituting equillibrium into the jacobian.
	\item Is it stable or unstable?
	\item The second part of this analysis will explain what happens when we add back the diffusion term.
	\item Does the system become unstable? If it does we do a bifurcation analysis to find which parameters $F$ and $k$ will result in turning patterns.
\end{itemize}
\subsection*{Numerical Simulation of the Gray Scott Equation}
% Describe why you're doing a numerical simulation
This subsection will document my attempt at a numerical simulation of the Gray Model. The numerical simulation is done using a finite difference approximation using 5 point stencil.The code for the simulation is in the appendix.
% Explain the Methodology
I use the method used by pearson in his paper on Complex patterns in a simple system.
% Figures showing the results of the simulation 
I need to complete this section
% Discussion of the results 
Then I will discuss the results and implications on what i just did.









